\documentclass[11pt]{article}

\usepackage[margin=1in]{geometry}
\usepackage{amsmath,amssymb,amsthm}
\usepackage{booktabs}
\usepackage{enumitem}
\usepackage{hyperref}
\usepackage{microtype}

\hypersetup{
  colorlinks=true,
  linkcolor=blue,
  citecolor=blue,
  urlcolor=blue
}
\setlength{\emergencystretch}{3em}

\newtheorem{assumption}{Assumption}
\newtheorem{proposition}{Proposition}
\newtheorem{definition}{Definition}

\newcommand{\stl}{\textsc{STL}}
\newcommand{\hal}{\mathrm{Hal}}
\newcommand{\baku}{\mathrm{Baku}}
\newcommand{\states}{\mathcal{S}}
\newcommand{\actions}{\mathcal{A}}
\newcommand{\expect}{\mathbb{E}}
\newcommand{\prob}{\mathbb{P}}
\newcommand{\valuefn}{V}
\newcommand{\policy}{\pi}

\title{A Search-and-Learning Solver for the \stl{} Handkerchief Game}
\author{Repository white-paper draft}
\date{May 28, 2026}

\begin{document}
\maketitle

\begin{abstract}
This draft documents the current solver architecture in the repository and
states the mathematical contract that must hold for the solver to be called
defensible. The implemented system models the handkerchief game as a
finite, stochastic, zero-sum, perfect-information game with simultaneous
move matrix games at each half-round. The code currently contains an
engine-derived exact finite-horizon minimax solver, a selective exact-second
search layer, tablebase pins, matrix-game Monte Carlo tree search (MCTS), a
neural value-and-policy network, supervised target generation, an
AlphaZero-style bootstrap loop, calibration gates, tournament evaluation,
and audit-pack checks. The latest documented run, Phase G iteration 3,
reaches held-out overall MSE 0.05934, a 45 percent improvement over the
pre-Phase-G baseline. The remaining rigor gap is not the existence of a
formal game value, but the approximation chain: finite horizon, selective
candidate sets, sampled MCTS values, and learned value-network error.
\end{abstract}

\section{Scope and Repository Snapshot}

The repository has shifted from a playable baseline toward a
Stockfish/AlphaZero-style solver. The central design choice is to keep the
canonical game in a fixed public-state legality regime rather than introduce
Bayesian type inference. In the current code, the strategic asymmetry is
encoded directly in legality: Baku may drop on second 61 during the leap
window, but Hal never may drop or check on second 61. This makes the solver
a perfect-information stochastic-game solver with simultaneous actions, not
a poker-style imperfect-information solver.

The last few weeks of commits can be summarized as follows.

\begin{itemize}[leftmargin=*]
  \item April 28 built the rigorous core: exact public states, LP minimax,
  engine-derived chance expansion, a firewall separating rigorous CFR code
  from shaping and bucketed abstractions, tablebase targets, selective
  exact-second search, and MCTS slices through selection, expansion, backup,
  search, transposition, and principal-line diagnostics.
  \item April 29--30 added the first value-target generation, the training
  harness, calibration metrics, tournament evaluation, subgame resolve, and
  the first bootstrap-loop API.
  \item May 18 hardened the bootstrap loop into a defensibility gate:
  stronger evaluator interfaces, audit packs, subgame-resolve hooks,
  source-weighted training, held-out calibration, and tests for legality,
  self-play, tournaments, and target generation.
  \item May 19--21 fixed corpus death-axis handling and CPR scaling, added
  corpus diagnostics, and expanded the pinned tablebase registry.
  \item May 24 added per-source calibration thresholds, a hidden-width
  architecture knob, a bounded LRU cache for exact LP recursion, and the
  Phase G wider LP-corpus script.
  \item May 26--27 recorded the Phase G + I-1 milestone and two subsequent
  bootstrap iterations. Iteration 3 converged at held-out MSE 0.05934.
\end{itemize}

The most important implementation files are:

\begin{center}
\begin{tabular}{ll}
\toprule
Path & Role \\
\midrule
\path{src/Game.py} & Authoritative half-round state machine \\
\path{src/Player.py}, \path{src/Referee.py} & Player and survival dynamics \\
\path{environment/legal_actions.py} & Actor-aware action legality \\
\path{environment/cfr/exact.py} & Exact public state, LP minimax, finite-horizon recursion \\
\path{environment/cfr/selective.py} & Candidate generation and selective minimax search \\
\path{environment/cfr/mcts.py} & Matrix-game-at-node MCTS \\
\path{environment/cfr/evaluator.py} & Terminal, tablebase, and value-net evaluators \\
\path{environment/cfr/tablebase.py} & Pinned tactical registry \\
\path{training/value_targets.py} & Exact and MCTS target generation \\
\path{training/train_value_net.py} & Supervised value/policy training \\
\path{training/bootstrap_loop.py} & Generation orchestration and hard gates \\
\path{training/calibration.py} & Held-out MSE, Brier score, reliability bins \\
\path{training/audit_pack.py} & Tactical audit records and drift gates \\
\path{hal/value_net.py} & Current neural architecture and 23-feature extractor \\
\bottomrule
\end{tabular}
\end{center}

\section{Canonical Game Model}

\subsection{Public State}

Let the two players be \(i \in \{\hal,\baku\}\). A public state
\(s \in \states\) is represented in the engine by the tuple
\[
s =
  (x_\hal,x_\baku,\tau_\hal,\tau_\baku,d_\hal,d_\baku,
   a_\hal,a_\baku,k,t,h,r,f,g,w),
\]
where \(x_i\) is the current cylinder time, \(\tau_i\) is total prior time
dead, \(d_i\) is the number of prior deaths, \(a_i\) records life status,
\(k\) is referee CPR count, \(t\) is game-clock seconds since 8:00:00 AM,
\(h \in \{1,2\}\) is the half-round, \(r\) is the round index, \(f\) is the
first dropper, \(g\) records terminal status, and \(w\) records the winner
when terminal. In code this is \path{environment/cfr/exact.py}'s
\path{ExactPublicState}.

The canonical opening state uses \(t=720\), i.e. 8:12:00 AM. The leap
window is \([3540,3600]\) in engine seconds, corresponding to the 8:59
minute and the inserted second 8:59:60.

\subsection{Half-Round Actions}

Each half-round has a dropper \(D(s)\) and checker \(C(s)\). The action is a
simultaneous pair
\[
a=(d,c), \qquad d \in A_D(s), \quad c \in A_C(s),
\]
where \(d\) is the drop second and \(c\) is the check second. Normal turns
have duration 60. Leap-window turns have duration 61.

The current actor-aware legality rule is:
\[
\begin{aligned}
  A_D(s) &= \{1,\ldots,60\} && \text{if } D(s)=\hal,\\
  A_D(s) &= \{1,\ldots,T(s)\} && \text{if } D(s)=\baku,\\
  A_C(s) &= \{1,\ldots,60\} && \text{for both players,}
\end{aligned}
\]
where \(T(s)=61\) only in the leap window and \(T(s)=60\) otherwise. This is
the decisive model reduction: Hal cannot act on second 61, and Baku's only
leap-second privilege is as dropper.

\subsection{Deterministic Resolution Before Death Chance}

For action \(a=(d,c)\), the checker succeeds iff \(c \ge d\). If the checker
succeeds, squandered time is
\[
\Delta x_C = \max(1,c-d).
\]
If the check fails, the checker receives the fixed failed-check penalty
\[
\Delta x_C = 60.
\]
If the checker cylinder reaches or exceeds 300 seconds, or if the checker
failed, the checker enters a death episode. The death duration is
\[
z = \min(x_C+\Delta x_C,300).
\]
If no death occurs, the transition is deterministic. If a death occurs, the
transition branches on survival.

\subsection{Survival Probability}

The referee computes survival probability as
\[
p_{\mathrm{survive}}(i,z,s)
 =
\left[1-\left(\frac{z}{300}\right)^3\right]_+
\cdot 0.85^{\tau_i/60}
\cdot \max(0.4,0.88^k)
\cdot \rho_i,
\]
with \(\rho_\hal=1.0\) and \(\rho_\baku=0.94\). The engine returns
probability zero at \(z \ge 300\). The chance branch is therefore
\[
\omega \in \{\mathrm{survive},\mathrm{die}\}, \qquad
\prob(\omega=\mathrm{survive})=p_{\mathrm{survive}}.
\]

\subsection{Utility}

All rigorous solver modules use terminal-only utility from Hal's
perspective:
\[
u_\hal(s) =
\begin{cases}
  +1, & \text{Hal is the terminal winner},\\
  -1, & \text{Baku is the terminal winner},\\
  0, & \text{draw or neutral terminal state}.
\end{cases}
\]
Nonterminal frontier states are not shaped in the exact solver. They are
recorded as unresolved probability mass when the finite horizon is exhausted.

\section{Game-Theoretic Foundation}

\begin{assumption}[Fixed engine model]
The Python engine is treated as the rules oracle. The solver is correct only
relative to this engine's transition and legality functions.
\end{assumption}

\begin{definition}[Finite-horizon value]
Let \(\valuefn_h(s)\) be Hal's value at state \(s\) with \(h\) half-rounds
remaining in the exact recursion. If \(s\) is terminal, then
\(\valuefn_h(s)=u_\hal(s)\). If \(h=0\) and \(s\) is nonterminal, then
\(\valuefn_0(s)=0\) and the unresolved mass is 1.
\end{definition}

For nonterminal \(s\) and \(h>0\), form a Hal-payoff matrix \(M_h(s)\) over
the legal exact seconds:
\[
M_h(s)_{d,c}
 =
\expect_{\omega \sim P(\cdot \mid s,d,c)}
\left[
  \valuefn_{h-1}(T(s,d,c,\omega))
\right].
\]
If Hal is the dropper, the row player maximizes \(M_h\); if Hal is the
checker, the signs are transposed so the row-player LP is still a maximin
solve.

The implemented row-player LP is the standard finite zero-sum matrix-game
program:
\[
\begin{aligned}
\max_{\sigma,v}\quad & v \\
\mathrm{s.t.}\quad & \sum_i \sigma_i M_{ij} \ge v
  \quad \forall j,\\
& \sum_i \sigma_i = 1,\\
& \sigma_i \ge 0.
\end{aligned}
\]
This appears as \path{solve_minimax} in \path{environment/cfr/exact.py}.

\begin{proposition}[Exact finite-horizon equilibrium value]
For a fixed engine state \(s\), fixed horizon \(h\), and full legal action
sets \(A_D(s),A_C(s)\), \path{solve_exact_finite_horizon} returns the value
of the finite-horizon stochastic zero-sum game induced by the engine and the
terminal-only utility.
\end{proposition}

\begin{proof}
The base case is terminal utility or an explicit unresolved frontier value.
For \(h>0\), every joint action is expanded through the engine. If death is
possible, the two chance branches are weighted by the referee probability;
otherwise the transition has probability one. The resulting \(M_h(s)\) is a
finite zero-sum matrix game. By von Neumann's minimax theorem, the LP value
equals the equilibrium value of that simultaneous move. Induction on \(h\)
then gives the finite-horizon stochastic-game value.
\end{proof}

The proposition is intentionally narrow. It proves the exact finite-horizon
solver, not the learned solver. The learned solver inherits strength only
through approximation quality and validation.

\section{Exact Core and Firewall}

The rigorous namespace \path{environment/cfr/} is guarded by tests that
prevent bucketed abstractions, reward shaping, observations, and route labels
from leaking into exact search. This is important: once shaped reward or a
bucket representative enters the exact core, the LP value is no longer a
statement about the canonical engine game.

The exact core currently provides:

\begin{itemize}[leftmargin=*]
  \item \path{ExactPublicState}: a full-precision public state key for cache
  and tablebase comparisons.
  \item \path{ExactGameSnapshot}: mutation-safe snapshot/restore around the
  authoritative engine.
  \item \path{enumerate_joint_actions}: full exact-second action enumeration.
  \item \path{expand_joint_action}: engine-derived chance expansion.
  \item \path{solve_exact_finite_horizon}: recursive exact minimax with
  unresolved frontier accounting.
  \item A bounded LRU cache keyed by public state, horizon, and perspective.
\end{itemize}

The latest cache change matters operationally. Phase G's wide LP grid would
otherwise duplicate deep recursion caches across worker processes and exhaust
memory. The current default cap is 30,000 entries per worker, controlled by
\path{STL_SOLVE_CACHE_MAXSIZE}.

\section{Selective Search and Tablebase}

Full-width exact search is expensive: normal turns have up to \(60 \times 60\)
joint actions and leap turns can include a 61st drop action for Baku. At
horizon 3 the tree becomes too large for broad corpus generation.

\path{environment/cfr/selective.py} defines a candidate generator that keeps:
\[
\{1,2,58,59,60,61\}
\]
where legal, plus local diagonals, safe-check boundaries, and overflow
threshold boundaries. The central candidate quantities are
\[
\mathrm{overflow}(x_C)=300-\lfloor x_C \rfloor
\]
and
\[
\mathrm{safe}(x_C)=299-\lfloor x_C \rfloor.
\]
Selective search solves the same minimax recursion over this reduced
candidate set and can audit against full-width exact search on small states.

The tablebase registry in \path{environment/cfr/tablebase.py} stores tactical
scenarios with pinned expected values or relational properties. Pinned entries
are exact ground-truth anchors for training and calibration. Phase F expanded
this registry from a tiny set to a broader suite of forced-overflow, clock,
leap-window, fatigue, TTD, asymmetric-death, and double-overflow cases.

\section{Matrix-Game MCTS}

The MCTS layer is not a vanilla single-agent UCT tree. Each node is itself a
simultaneous-move matrix game. A node stores candidate drop seconds, candidate
check seconds, a value matrix \(Q\), visit counts \(N\), a prior matrix \(P\),
and transposition links.

At selection time it forms
\[
\widetilde{Q}_{ij}
 =
Q_{ij}
 +
c_{\mathrm{puct}}
P_{ij}
\frac{\sqrt{N_{\mathrm{node}}}}{1+N_{ij}},
\]
then solves the matrix game over \(\widetilde{Q}\) and samples a joint action
from the resulting dropper/checker strategies. The engine is stepped into the
child; if a death is possible, MCTS samples the survival branch. At an
unexpanded nonterminal leaf, the supplied evaluator returns a
Hal-perspective value and optional dropper/checker policy priors. Backup uses
a running mean and does not flip signs, because every value is already in
Hal's perspective.

The current MCTS implementation includes:

\begin{itemize}[leftmargin=*]
  \item evaluator-supplied policy priors over 61 dropper and 61 checker
  seconds;
  \item transposition reuse keyed by exact public state;
  \item principal-line extraction via most-visited cells;
  \item optional root subgame re-solving at critical states.
\end{itemize}

\section{Value Network}

The current neural evaluator in \path{hal/value_net.py} uses 23 scalar
features:
normalized cylinders, total time dead, death counts, safe budgets, projected
fail-survival probabilities, role bit, normalized clock, round/half
features, LSR variation one-hot features, leap-turn bit, rounds until the
leap window, referee CPR count, active LSR bit, and leap proximity.

The network is:
\[
\phi_\theta(x)
 =
\mathrm{ReLU}(W_2 \mathrm{ReLU}(W_1x+b_1)+b_2),
\]
with a scalar value head
\[
\hat v_\theta(x)=\tanh(w_v^\top \phi_\theta(x)+b_v)
\]
and a policy head with 122 logits: 61 for dropper seconds and 61 for checker
seconds. The historical default is hidden width 64. Phase I-1 uses hidden
width 128, giving about 35.5K parameters, up from about 13.7K.

Training minimizes source-weighted value MSE plus masked policy cross-entropy:
\[
\mathcal{L}(\theta)
 =
\frac{1}{n}
\sum_{m=1}^n
\alpha_{s_m}
(\hat v_\theta(x_m)-y_m)^2
+
\lambda
\left[
  H(\pi^D_m,\hat\pi^D_\theta)
  +H(\pi^C_m,\hat\pi^C_\theta)
\right],
\]
where illegal seconds are masked before the softmax. Source weights are used
because terminal and tablebase anchors can be numerically rare but
strategically decisive.

\section{Target Generation and Bootstrap Loop}

The repository no longer treats horizon-1 exact labels as sufficient.
Horizon-1 labels are too often unresolved and therefore blind to leap-second
routing. The current target hierarchy is:

\begin{enumerate}[leftmargin=*]
  \item terminal state: terminal utility, source \path{terminal};
  \item pinned tablebase state: exact pin, source \path{tablebase};
  \item LSR-significant state: exact horizon 3, source
  \path{exact_horizon_3};
  \item active-LSR or near-overflow state: exact horizon 2, source
  \path{exact_horizon_2};
  \item otherwise: excluded from the exact corpus and left to MCTS bootstrap.
\end{enumerate}

The AlphaZero-style loop is:

\begin{enumerate}[leftmargin=*]
  \item train a value/policy network on exact and tablebase anchors;
  \item run MCTS using that network at leaves;
  \item record root MCTS values and root policies as bootstrap labels;
  \item merge exact anchors back in to prevent self-referential drift;
  \item train the next network;
  \item reject the generation if calibration gates fail.
\end{enumerate}

This follows the expert-iteration pattern used by AlphaZero and related
search-and-learning systems: search improves the policy/value estimate, then
the network distills search output into a cheaper evaluator.

\section{Current Empirical Results}

The latest documented chain is Phase G + I-1 through Phase G iteration 3.

\begin{center}
\begin{tabular}{lrr}
\toprule
Milestone & Held-out overall MSE & Comment \\
\midrule
gen-2 v2 baseline & 0.10841 & Pre-Phase-G best \\
gen-3 v2 & 0.132 & Plateau/regression \\
Phase G + I-1 & 0.070 & Strict per-source gate passed \\
Phase G iter 2 & 0.06000 & Policy-improvement signal remains \\
Phase G iter 3 & 0.05934 & Converged under current data/architecture \\
\bottomrule
\end{tabular}
\end{center}

The Phase G LP corpus had 7,168 candidate states and emitted 4,812 records:
2,938 exact-horizon-3, 1,871 exact-horizon-2, and 3 direct tablebase records.
After anchors and bootstrap merge, iteration 2 and iteration 3 trained on
6,357-record corpora with 864 MCTS-bootstrap labels plus exact/tablebase and
terminal anchors.

The strict held-out gate checks per-source MSE, tablebase MSE, required
source presence, and unresolved probability. Phase G iteration 3 passed with:

\begin{center}
\begin{tabular}{lrr}
\toprule
Source & n & MSE \\
\midrule
terminal & 24 & 0.498 \\
tablebase & 19 & 0.007 \\
exact\_horizon\_2 & 128 & 0.012 \\
exact\_horizon\_3 & 64 & 0.004 \\
\bottomrule
\end{tabular}
\end{center}

This is strong evidence that the supervised loop is now data-limited under
the current target distribution. It is not yet proof that the root game is
solved.

\section{Relationship to AlphaZero, CFR, and Poker Solvers}

The architecture resembles AlphaZero in its search-distillation loop:
MCTS guided by a neural value/policy network produces improved targets, and
the next generation network learns those targets. The difference is that
each \stl{} node is a simultaneous zero-sum matrix game rather than a
single-player choice by the side to move.

The architecture resembles CFR systems in its use of equilibrium values for
local matrix games and its concern with exploitability, but the current core
does not perform regret minimization over information sets. The namespace
name \path{environment/cfr/} is therefore historical and partially
aspirational. The rigorous exact solver is better described as finite-horizon
stochastic minimax with LP solves at simultaneous-move nodes.

DeepStack, Libratus, Pluribus, Deep CFR, and ReBeL solve harder
imperfect-information poker settings by combining abstraction, regret
minimization, depth-limited solving, and learned value functions. Those tools
are relevant literature, but they are not required if the canonical \stl{}
model remains fully public after the legality reduction. If future canon
work reintroduces hidden memory or belief states, the correct next framework
would be public-belief-state solving or CFR over information sets. Under the
current repo contract, that would add complexity without improving the
formal model being solved.

\section{What Is Rigorous and What Is Approximate}

The mathematically rigorous pieces are:

\begin{itemize}[leftmargin=*]
  \item the engine-derived transition model, assuming the engine is the
  accepted rule oracle;
  \item the exact finite-horizon full-width solver;
  \item LP minimax at each simultaneous move;
  \item terminal-only utility in the rigorous core;
  \item pinned tablebase scenarios whose expected values are exact by
  construction.
\end{itemize}

The approximate pieces are:

\begin{itemize}[leftmargin=*]
  \item finite-horizon cutoff and unresolved mass;
  \item selective candidate generation;
  \item MCTS sampling variance and finite iteration budget;
  \item learned value-network approximation error;
  \item policy-head targets distilled from exact or MCTS root strategies;
  \item empirical calibration gates, which measure but do not prove global
  exploitability.
\end{itemize}

The most important distinction for the white paper is this: a low held-out
MSE says the evaluator predicts the current target corpus well. It does not
by itself bound exploitability across the full game tree. To make the solver
paper-grade, the next version should add explicit exploitability and
candidate-set regret audits on increasingly broad state samples.

\section{Best Next Steps}

\subsection{1. Expand the Held-Out and Training Anchors}

Phase G iteration 3 appears to have reached a data ceiling. The highest-value
next step is Phase F-2: expand the tablebase and held-out registry beyond the
current forced-overflow-heavy family. Add pins around:

\begin{itemize}[leftmargin=*]
  \item non-299 near-overflow thresholds;
  \item non-leap midgame states where the correct play routes into LSR
  variation 2 several rounds later;
  \item asymmetric TTD/death-count states that are not already collapsed by
  the 23-feature representation;
  \item role-swapped leap-window states;
  \item states where Baku's second-61 drop is legal but not trivially winning.
\end{itemize}

Do this before increasing model size again. The hidden-128 run already showed
that larger capacity without enough diverse data overfits.

\subsection{2. Add High-Variance MCTS Reanalysis}

Phase H should select states where MCTS value has high seed variance, high
policy entropy, or disagreement with exact/selective labels, then re-run
those states at higher iteration counts. This is a better use of compute than
uniformly enlarging the bootstrap grid.

\subsection{3. Quantify Exploitability}

The repo has diagnostics for exact matrix games, but the white paper needs a
headline exploitability story for the learned agent. Add:

\begin{itemize}[leftmargin=*]
  \item exact best-response gaps on small full-width subgames;
  \item selective-vs-full-width candidate regret on audited states;
  \item MCTS seed confidence intervals at fixed roots;
  \item value error stratified by state feature axes, not only source class.
\end{itemize}

\subsection{4. Promote Policy Validation}

The value head is currently better justified than the policy head. Add policy
metrics:

\begin{itemize}[leftmargin=*]
  \item root-action KL against exact/selective strategies on held-out states;
  \item illegal-mass checks after masking;
  \item top-action agreement where the exact equilibrium has low entropy;
  \item entropy calibration for mixed strategies.
\end{itemize}

\subsection{5. Run Canonical Principal-Line Validation}

\path{scripts/validate_against_manga.py} currently supports the intended
validation path, but its default evaluator is still the terminal-only
placeholder. The next useful milestone is to load
\path{checkpoints/gen_phaseG_iter3/best.pt}, run the canonical R1T1-to-R9T2
scenario against \path{BakuLSREngineeringTeacher}, and report the principal
line, root values, and deviations from the documented canonical strategy.

\subsection{6. Decide Whether This Is a Solver or a Strong Agent}

For a paper, terminology matters. The current system is defensibly a
calibrated search-and-learning agent with exact local solvers. It is not yet
a full-game solved equilibrium strategy unless the paper can bound
exploitability over the relevant starting distribution. The next writing
standard should be:

\begin{itemize}[leftmargin=*]
  \item call exact finite-horizon outputs ``solves'';
  \item call Phase G iter 3 a ``calibrated evaluator'';
  \item call MCTS plus evaluator a ``search policy'';
  \item reserve ``solver'' for the packaged system only after exploitability
  and candidate-regret bounds are included.
\end{itemize}

\section{Conclusion}

The repository now has the right skeleton for a rigorous solver paper:
canonical rules, a formal public-state model, engine-derived stochastic
transitions, exact finite-horizon minimax, tablebase pins, matrix-game MCTS,
value/policy distillation, calibration gates, and reproducible scripts. The
main scientific work remaining is to turn empirical calibration into a
stronger approximation argument: broader anchors, high-variance reanalysis,
explicit exploitability estimates, and principal-line validation on the
canonical scenario.

\begin{thebibliography}{99}

\bibitem{vonneumann1928}
John von Neumann.
\newblock Zur Theorie der Gesellschaftsspiele.
\newblock \emph{Mathematische Annalen}, 100:295--320, 1928.
\newblock \url{https://doi.org/10.1007/BF01448847}

\bibitem{shapley1953}
Lloyd S. Shapley.
\newblock Stochastic games.
\newblock \emph{Proceedings of the National Academy of Sciences}, 39(10):1095--1100, 1953.
\newblock \url{https://doi.org/10.1073/pnas.39.10.1095}

\bibitem{zinkevich2007}
Martin Zinkevich, Michael Johanson, Michael Bowling, and Carmelo Piccione.
\newblock Regret minimization in games with incomplete information.
\newblock \emph{Advances in Neural Information Processing Systems}, 2007.
\newblock \url{https://papers.nips.cc/paper_files/paper/2007/hash/08d98638c6fcd194a4b1e6992063e944-Abstract.html}

\bibitem{browne2012}
Cameron B. Browne, Edward Powley, Daniel Whitehouse, Simon M. Lucas, Peter I. Cowling, Philipp Rohlfshagen, Stephen Tavener, Diego Perez, Spyridon Samothrakis, and Simon Colton.
\newblock A survey of Monte Carlo tree search methods.
\newblock \emph{IEEE Transactions on Computational Intelligence and AI in Games}, 4(1):1--43, 2012.
\newblock \url{https://doi.org/10.1109/TCIAIG.2012.2186810}

\bibitem{deepstack2017}
Matej Moravcik, Martin Schmid, Neil Burch, Viliam Lisy, Dustin Morrill, Nolan Bard, Trevor Davis, Kevin Waugh, Michael Johanson, and Michael Bowling.
\newblock DeepStack: Expert-level artificial intelligence in heads-up no-limit poker.
\newblock \emph{Science}, 356(6337):508--513, 2017.
\newblock \url{https://doi.org/10.1126/science.aam6960}

\bibitem{alphazero2018}
David Silver, Thomas Hubert, Julian Schrittwieser, Ioannis Antonoglou, Matthew Lai, Arthur Guez, Marc Lanctot, Laurent Sifre, Dharshan Kumaran, Thore Graepel, Timothy Lillicrap, Karen Simonyan, and Demis Hassabis.
\newblock A general reinforcement learning algorithm that masters chess, shogi, and Go through self-play.
\newblock \emph{Science}, 362(6419):1140--1144, 2018.
\newblock \url{https://doi.org/10.1126/science.aar6404}

\bibitem{libratus2018}
Noam Brown and Tuomas Sandholm.
\newblock Superhuman AI for heads-up no-limit poker: Libratus beats top professionals.
\newblock \emph{Science}, 359(6374):418--424, 2018.
\newblock \url{https://doi.org/10.1126/science.aao1733}

\bibitem{pluribus2019}
Noam Brown and Tuomas Sandholm.
\newblock Superhuman AI for multiplayer poker.
\newblock \emph{Science}, 365(6456):885--890, 2019.
\newblock \url{https://doi.org/10.1126/science.aay2400}

\bibitem{deepcfr2019}
Noam Brown, Adam Lerer, Sam Gross, and Tuomas Sandholm.
\newblock Deep counterfactual regret minimization.
\newblock \emph{Proceedings of the 36th International Conference on Machine Learning}, 2019.
\newblock \url{https://proceedings.mlr.press/v97/brown19b.html}

\bibitem{rebel2020}
Noam Brown, Anton Bakhtin, Adam Lerer, and Qucheng Gong.
\newblock Combining deep reinforcement learning and search for imperfect-information games.
\newblock \emph{Advances in Neural Information Processing Systems}, 2020.
\newblock \url{https://proceedings.neurips.cc/paper/2020/hash/c61f571dbd2fb949d3fe5ae1608dd48b-Abstract.html}

\end{thebibliography}

\end{document}
